% Part: first-order-logic
% Chapter: introduction
% Section: semantic-notions

\documentclass[../../../include/open-logic-section]{subfiles}

\begin{document}

\olfileid{fol}{int}{sem}

\olsection{مفاهیم معنایی}

پیش‌تر گفتیم هنگامی که برقرار بودن $\Sat{M}{!A}[s]$ را بررسی می‌کنیم،
(برای سهولت) می‌گذاریم $s$ به همهٔ !!{variable}s مقدار دهد، اما فقط
مقادیر نسبت‌داده‌شده به !!{variable}s موجود در~$!A$ به کار می‌روند. در
واقع، تنها مقادیر متغیرهای \emph{آزاد} در~$!A$ اهمیت دارند. البته چون
دقیق عمل می‌کنیم، این واقعیت را اثبات خواهیم کرد. از آنجا که
!!{sentence}s متغیر آزاد ندارند، در ارضاشدن یا نشدن آن‌ها در
!!a{structure}،~$s$ هیچ اهمیتی ندارد. بنابراین، هنگامی که $!A$
!!a{sentence} است، می‌توانیم $\Sat{M}{!A}$ را به معنای
«$\Sat{M}{!A}[s]$ برای هر~$s$» تعریف کنیم؛ و اتفاقاً این شرط درست است
هرگاه و تنها هرگاه $\Sat{M}{!A}[s]$ برای دست‌کم یک~$s$ درست باشد. برای
به‌دست‌دادن تعریفی کارآمد از ارضای !!{formula}s باید گمارش‌های
!!{variable} را معرفی کنیم، اما ارضای !!{sentence}s از این گمارش‌ها
مستقل است.

پس از تعریف «$\Sat{M}{!A}$»، می‌دانیم تا آنجا که به !!{sentence}s
منطق مرتبهٔ اول مربوط است، «حالت» و «صادق در» چه معنایی دارند. سپس بر
پایهٔ تعریف $\Sat{M}{!A}$ برای !!{sentence}s می‌توانیم مفاهیم معنایی
پایهٔ اعتبار، استلزام و ارضاپذیری را تعریف کنیم. یک جمله معتبر است،
$\Entails !A$، اگر هر !!{structure} آن را ارضا کند. مجموعه‌ای از
!!{sentence}s آن را نتیجه می‌دهد، $\Gamma \Entails !A$، اگر هر
!!{structure} که همهٔ !!{sentence}s موجود در~$\Gamma$ را ارضا می‌کند،
~$!A$ را نیز ارضا کند. و مجموعه‌ای از !!{sentence}s ارضاپذیر است اگر
!!{structure}ی وجود داشته باشد که همهٔ !!{sentence}s آن را هم‌زمان
ارضا کند.

چون !!{formula}s به‌طور استقرایی تعریف شده‌اند و ارضا نیز به‌نوبهٔ خود
با استقرا بر ساختار !!{formula}s تعریف شده است، می‌توانیم با استقرا
خواص معناشناسی خود را اثبات و مفاهیم معنایی تعریف‌شده را به هم مرتبط
کنیم. برخی از این خواص را گرد می‌آوریم و اثبات می‌کنیم؛ تا اندازه‌ای
چون هر یک به‌تنهایی جالب‌اند، اما بیشتر از آن رو که بسیاری از آن‌ها
هنگام بررسی رابطهٔ معناشناسی و دستگاه‌های !!{derivation} سودمند خواهند
بود. برای این کار، باید برخی مفاهیم و عملیات نحوی را که هنوز ذکر
نکرده‌ایم نیز (به‌دقت، یعنی با استقرا) تعریف کنیم.

\end{document}
